\documentclass[12pt]{article}
\usepackage{amsmath,amsfonts,fullpage,amsthm,graphicx}

%Style section
%\setlength{\textheight}{10in}
%\setlength{\textwidth}{7in}
%\hoffset=-0.75in
\pagestyle{plain}
% \setlength{\topmargin}{0in}
% \setlength{\headsep}{0in}
% \setlength{\oddsidemargin}{0in}
% \setlength{\evensidemargin}{0in}

\renewcommand{\qedsymbol}{$\blacksquare$}


\newcommand{\emptybox}[2][\textwidth]{%
	\begingroup
	\setlength{\fboxsep}{-\fboxrule}%
	\noindent\framebox[#1]{\rule{0pt}{#2}}%
	\endgroup
}



\makeatletter
\def\saveenum{\xdef\@savedenum{\the\c@enumi\relax}}
\def\resetenum{\global\c@enumi\@savedenum}
\makeatother


\theoremstyle{definition}
\newtheorem*{definition}{Definition}

\newtheorem*{theorem}{Theorem}

\theoremstyle{remark}
\newtheorem*{remark}{Remark}

\newtheorem{example}{Example}

\newtheorem{lemma}{Lemma}



\begin{document}

\pagestyle{empty} %use this to get rid of page numbers

\noindent
{Math 166 \hskip 1 truein  Name:\ \hrulefill}

\noindent
%\vskip 0.3truein
%\bigskip
%\noindent
{All series test practice}
\noindent \hspace{2.5in} Section Number: \hrulefill
\bigskip


%\textit{Quick Review:}\\
%\textbf{Alternating Series Test:} If the alternating series \begin{equation*}
 %   \sum\limits_{n=1}^\infty (-1)^{n+1}a_n = a_1 - a_2 + a_3 - a_4 + a_5 - a_6 + \ldots \;\;\;\; \text{where }a_n \geq 0 \text{ for all n}
%\end{equation*}
%satisfies the conditions: \begin{itemize}
 %   \item $a_{n+1} \leq a_n$ for all $n$ \\
 %  \item $\lim_{n\rightarrow \infty} a_n = 0$
% \end{itemize}
% Then the series converges.
% \vspace{.2in}
% \hrule \vspace{.2in}
Determine whether the series converges conditionally, converges absolutely, or diverges.
\begin{enumerate}
	 \item 
	\begin{equation*}
		\sum\limits_{n=2}^\infty \frac{(-1)^n}{\ln(n)} 
	\end{equation*}

\begin{center}
	\raisebox{-.3\height}{\emptybox[.25in]{.25in}} Diverges \hspace{.5 in} \raisebox{-.3\height}{\emptybox[.25in]{.25in}} Converges conditionally \hspace{.5 in} \raisebox{-.3\height}{\emptybox[.25in]{.25in}} Converges absolutely
\end{center}  
Reason:
\vfill
\vfill
	

    \item 
    \begin{equation*}
        \sum\limits_{n=0}^\infty \cos(\pi n)\frac{n+1}{n+3} 
    \end{equation*}
\begin{center}
	\raisebox{-.3\height}{\emptybox[.25in]{.25in}} Diverges \hspace{.5 in} \raisebox{-.3\height}{\emptybox[.25in]{.25in}} Converges conditionally \hspace{.5 in} \raisebox{-.3\height}{\emptybox[.25in]{.25in}} Converges absolutely
\end{center}  
Reason:
\vfill
\vfill
\newpage
%\textbf{The Ratio Test:} \begin{itemize}
  %  \item If $\lim_{n\rightarrow \infty} \left| \frac{a_{n+1}}{a_n}\right| < 1$, then the series $\sum\limits_{n=1}^\infty a_n$ is absolutely convergent.
   % \item If $\lim_{n\rightarrow \infty} \left| \frac{a_{n+1}}{a_n}\right| > 1$, then the series is divergent.
    %\item If $\lim_{n\rightarrow \infty} \left| \frac{a_{n+1}}{a_n}\right| = 1$, then the ratio test is \textit{inconclusive}.
% \end{itemize}
\vspace{.2in}
 \vspace{.2in}


     \item  \begin{equation*}
        \sum\limits_{n=1}^\infty \frac{n}{5^n}
    \end{equation*}
    \begin{center}
    	\raisebox{-.3\height}{\emptybox[.25in]{.25in}} Diverges \hspace{.5 in} \raisebox{-.3\height}{\emptybox[.25in]{.25in}} Converges conditionally \hspace{.5 in} \raisebox{-.3\height}{\emptybox[.25in]{.25in}} Converges absolutely
    \end{center}  
    Reason:
    \vfill
    \vfill
   
    \item   \begin{equation*}
        \sum\limits_{n=1}^\infty \frac{2^n\ln(n)}{n!}
    \end{equation*}
    \begin{center}
    	\raisebox{-.3\height}{\emptybox[.25in]{.25in}} Diverges \hspace{.5 in} \raisebox{-.3\height}{\emptybox[.25in]{.25in}} Converges conditionally \hspace{.5 in} \raisebox{-.3\height}{\emptybox[.25in]{.25in}} Converges absolutely
    \end{center}  
    Reason:
    \vfill
    \vfill

    \item \begin{equation*}
        \sum\limits_{n=1}^\infty \frac{n^n}{n!}
    \end{equation*}
    \begin{center}
    	\raisebox{-.3\height}{\emptybox[.25in]{.25in}} Diverges \hspace{.5 in} \raisebox{-.3\height}{\emptybox[.25in]{.25in}} Converges conditionally \hspace{.5 in} \raisebox{-.3\height}{\emptybox[.25in]{.25in}} Converges absolutely
    \end{center}  
    Reason:
    \vfill
    \vfill


\end{enumerate}
\end{document}